% ============================================================
% CAPITOLO 32 - RIEPILOGO DELLA PARTE VI
% ============================================================
\chapter{Riepilogo della Parte VI}
\label{chap:riepilogo_vi}

\section{Sintesi dei rivolti della \sigla{V}{7}{}}
\label{sec:sintesi_rivolti_v7}

\begin{center}
\begin{tabularx}{\linewidth}{lLLLLL}
\toprule
\textbf{Stato} & \textbf{Cifratura} & \textbf{Nota al basso} & \textbf{Risolve su} & \textbf{Basso obbligato} & \textbf{Dettaglio} \\
\midrule
Fond. & 7 & fondamentale & I & no & Sez. \ref{subsec:v7_fondamentale} \\
I rv. & 6/5 & sensibile & I & sì (sale) & Sez. \ref{subsec:v65} \\
II rv. & 4/3 & quinta & I (o \sigla{I}{6}{}) & no & Sez. \ref{subsec:v43} \\
III rv. & 2 & settima & \sigla{I}{6}{} & sì (scende) & Sez. \ref{subsec:v2} \\
\bottomrule
\end{tabularx}
\end{center}

\section{\sigla{VII}{7}{} maggiore contro \sigla{II}{7}{} minore}
\label{sec:sintesi_vii_ii}

\begin{center}
\small
\begin{tabular}{lcc}
\toprule
\textbf{Proprietà} & \textbf{\sigla{VII}{7}{} (modo maggiore)} & \textbf{\sigla{II}{7}{} (modo minore)} \\
\midrule
Struttura & semidiminuita (m-m-M) & semidiminuita (m-m-M) \\
Funzione & dominante & sottodominante \\
Risoluzione naturale & tonica (I) & dominante (V) \\
Preparazione della settima & non necessaria & obbligatoria \\
Quinta diminuita a risoluzione obbligata & sì (scende) & sì (scende) \\
Rivolto principale & \sigla{VII}{6}{5} $\rightarrow$ \sigla{I}{6}{} & \sigla{II}{6}{5} $\rightarrow$ V \\
Note (es. Do M / La m) & Si-Re-Fa-La & Si-Re-Fa-La \\
\bottomrule
\end{tabular}
\end{center}

Dettaglio nel Capitolo \ref{chap:vii7_ii7}.

\section{Quadro riassuntivo delle specie}
\label{sec:quadro_specie}

\begin{center}
\small
\begin{tabularx}{\linewidth}{lLLLL}
\toprule
\textbf{Specie} & \textbf{Esempio} & \textbf{Risoluzione tipica} & \textbf{Moto della settima} & \textbf{Cap.} \\
\midrule
1\textsuperscript{a} dominante & \sigla{V}{7}{} $\rightarrow$ I & settima $\searrow$, sensibile $\nearrow$ & scende & \ref{chap:settima_dominante} \\
2\textsuperscript{a} minore & \sigla{II}{7}{} $\rightarrow$ V, \sigla{i}{7}{} $\rightarrow$ iv & quinta discendente & scende & \ref{chap:seconda_specie} \\
3\textsuperscript{a} semidiminuita & \sigla{VII}{7}{} $\rightarrow$ I, \sigla{II}{7}{} $\rightarrow$ V & tonica (sensibile $\nearrow$) o dominante & scende & \ref{chap:vii7_ii7} \\
4\textsuperscript{a} maggiore & \sigla{I}{7}{} $\rightarrow$ IV, \sigla{VI}{7}{} $\rightarrow$ II & quinta discendente & scende & \ref{chap:quarta_quinta_specie} \\
5\textsuperscript{a} diminuita & \sigla{VII}{7}{} $\rightarrow$ i & tonica, sensibile $\nearrow$ & scende & \ref{chap:quarta_quinta_specie} \\
6\textsuperscript{a} minore-maggiore & \sigla{i}{7}{} $\rightarrow$ iv & quinta discendente & \textbf{sale} (è la sensibile) & \ref{chap:specie_rare} \\
7\textsuperscript{a} con 5\textsuperscript{a} ecc. & \sigla{III}{7}{} $\rightarrow$ VI & 5\textsuperscript{a} ecc.\ $\nearrow$, settima $\searrow$ & scende (la 5\textsuperscript{a} ecc.\ sale) & \ref{chap:specie_rare} \\
\bottomrule
\end{tabularx}
\end{center}

Tutti gli esempi si trasportano in qualsiasi tonalità con la stessa condotta delle voci: cambiano soltanto i nomi delle note e le alterazioni in chiave.
